\section*{第 \ref{ch:b2:atomic-orbitals} 章——原子轨道：形状、能量与大小}

\begin{solution}{exo:b2:atomic-orbitals:1}
3p：$n = 3$，$l = 1$；有 1 个径向节点（$n - l - 1$）和 1 个角向节面。4d：$n = 4$，
$l = 2$；有 1 个径向节点和 2 个角向节面。
\end{solution}

\begin{solution}{exo:b2:atomic-orbitals:2}
$P_{2p} \propto r^4\mathrm e^{-r}$（$r$ 以 $a_0$ 为单位）：$d\ln P/dr = 4/r - 1 = 0$，得 $r = 4$：
$4a_0 = \qty{212}{pm}$。
\end{solution}

\begin{solution}{exo:b2:atomic-orbitals:3}
$1 - \mathrm e^{-2}(1 + 2 + 2) = 1 - 5\mathrm e^{-2} = 0.323$。
\end{solution}

\begin{solution}{exo:b2:atomic-orbitals:4}
四个瓣位于 $x$ 轴与 $y$ 轴之间，处于 $xy$ 平面内，符号交替变化
（$xy > 0$ 处为正）；节面为 $xz$ 面和 $yz$ 面。径向节点数：
$n - l - 1 = 0$。
\end{solution}

\begin{solution}{exo:b2:atomic-orbitals:5}
氧，$(1\mathrm s)^2(2\mathrm s, 2\mathrm p)^6$：$\sigma = 5 \times 0.35 + 2 \times 0.85
= 3.45$，$Z^* = 4.55$；半径 $4/4.55 = 0.88a_0$，\qty{47}{pm}。
\end{solution}

\begin{solution}{exo:b2:atomic-orbitals:6}
电离时首先失去束缚最弱的电子。4s 电子
（\qty{-14}{eV}）远比 3d 电子（\qty{-59}{eV}）束缚弱：
两个 4s 电子先离开，然后再失去一个 3d 电子，得到 $[\ce{Ar}]\,3\mathrm d^5$。
\end{solution}

\begin{solution}{exo:b2:atomic-orbitals:7}
锂，2s：$Z^* = 3 - 2 \times 0.85 = 1.30$，$I = 13.6 \times 1.30^2/4 =
\qty{5.75}{eV}$（实测 5.39，偏高 7~\%）。钠，3s：$Z^* = 11 - 8 \times 0.85 -
2 \times 1.00 = 2.20$，$I = 13.6 \times 2.20^2/9 = \qty{7.3}{eV}$（实测 5.14）：
斯莱特规则作为粗略模型，高估了最外层电子的束缚，
且沿族向下越来越严重。
\end{solution}

\begin{solution}{exo:b2:atomic-orbitals:8}
钠：$Z^* = 2.20$，半径 $9/2.20 = 4.1a_0$（\qty{216}{pm}）。氯，3p：
$\sigma = 6 \times 0.35 + 8 \times 0.85 + 2 \times 1.00 = 10.90$，$Z^* = 6.10$，半径
$9/6.10 = 1.5a_0$（\qty{78}{pm}）。两者位于同一电子层，但氯的最外层电子
所受的有效电荷几乎是钠的三倍，因而被拉得更近。
\end{solution}

\begin{solution}{exo:b2:atomic-orbitals:9}
两者的角向部分均为 $1/\sqrt{4\pi}$，因此积分化为
\[
  \int_0^\infty R_{1s}R_{2s}r^2\,dr \propto \int_0^\infty (2 - r)r^2\mathrm e^{-3r/2}\,dr =
  2 \times \frac{2!}{(3/2)^3} - \frac{3!}{(3/2)^4} = \frac{32}{27} - \frac{32}{27} = 0 .
\]
\end{solution}

\begin{solution}{exo:b2:atomic-orbitals:10}
$\int_0^\infty N^2r^2\mathrm e^{-r}r^2\,dr = N^2 \times 4! = 1$：$N = 1/\sqrt{24} =
1/(2\sqrt6)$，与表中一致。
\end{solution}

\begin{solution}{exo:b2:atomic-orbitals:11}
$\langle r\rangle = 1.5a_0$，最可几半径为 $a_0$。$P(r) = 4r^2\mathrm e^{-2r}$ 从零开始急剧上升
随后缓慢下降：$r$ 较大处的长尾对
平均值的影响大于对极大值位置的影响。
\end{solution}

\begin{solution}{exo:b2:atomic-orbitals:12}
$Y_{p_z} \propto z/r = \cos\theta$ 在 $\theta = \pi/2$ 处为零，即 $xy$ 平面。
$\psi$ 的符号与 $z$ 相同，上方为正，下方为负。在该平面上
概率密度为零。
\end{solution}

\begin{solution}{pb:b2:atomic-orbitals:1}
\textbf{1.} $\int|\psi|^2\,dV = \frac{4\pi}{\pi a_0^3}\int_0^\infty r^2\mathrm e^{-2r/a_0}\,dr =
\frac{4}{a_0^3} \times \frac{a_0^3}{4} = 1$。
\textbf{2.} $P(r) = 4\pi r^2|\psi|^2 = (4/a_0^3)\,r^2\mathrm e^{-2r/a_0}$。
\textbf{3.} $dP/dr = 0$：$2/r - 2/a_0 = 0$，$r = a_0$。
\textbf{4.} $\langle r\rangle = (4/a_0^3)\int_0^\infty r^3\mathrm e^{-2r/a_0}\,dr = (4/a_0^3)
\times 6(a_0/2)^4 = 1.5a_0$。
\textbf{5.} 分布是偏斜的，在 $r$ 较大处有长尾。
\textbf{6.} $|\psi|^2$ 是单位体积内的密度，在原子核处最大；$P(r)$
计入的是球壳体积 $4\pi r^2\,dr$，它在 $r = 0$ 处为零。
\textbf{7.} \qty{52.9}{pm}。
\textbf{8.} 令 $x = r/a_0$：$\int_{x_0}^\infty 4x^2\mathrm e^{-2x}\,dx$；两次分部积分得
$= \mathrm e^{-2x_0}(2x_0^2 + 2x_0 + 1)$。
\textbf{9.} $5\mathrm e^{-2} = 0.677$。
\textbf{10.} $13\mathrm e^{-4} = 0.238$。
\textbf{11.} $\mathrm e^{-2x}(1 + 2x + 2x^2) = 0.10$，$x \approx 2.66$：\qty{141}{pm}。
\textbf{12.} 密度在任何距离处都不为零，没有明确的边界。但 90~\% 的概率
位于 \qty{141}{pm} 之内：这是一个实用的大小。
\textbf{13.} 2s：有 1 个径向节点（半径为 $2a_0$ 的球面），无角向节面。2p：有 1 个
角向节面（一个平面），无径向节点。
\textbf{14.} 位于 $r = 2a_0$ 处。
\textbf{15.} $4a_0$（练习 2）。
\textbf{16.} $d\ln P/dr = 2/r - 2/(2 - r) - 1 = 0$ 给出 $r^2 - 6r + 4 = 0$，$r = 3 \pm \sqrt5$：
$0.76a_0$（较小的内极大值）和 $5.24a_0$（主极大值）。
\textbf{17.} 2s，通过其内极大值。在多电子原子中，2s 电子
穿透 1s 电子层，所受屏蔽小于 2p 电子，
能量更低。
\textbf{18.} 所有半径都除以 $Z = 6$：2p 为 $0.67a_0$（\qty{35}{pm}）；2s 的极大值
位于 $0.13a_0$ 和 $0.87a_0$ 处。
\textbf{19.} C 3.25，N 3.90，O 4.55。
\textbf{20.} $\varepsilon = -13.6\,Z^{*2}/4$：C \qty{-35.9}{eV}，N \qty{-51.7}{eV}，O
\qty{-70.4}{eV}。
\textbf{21.} 同本章所述：\qty{11.5}{eV}（实测 11.26）。
\textbf{22.} N：原子 $5 \times (-13.6 \times 3.90^2/4) = \qty{-258.6}{eV}$；\ce{N+}，$Z^* =
4.25$，$4 \times (-13.6 \times 4.25^2/4) = \qty{-245.7}{eV}$；$I = \qty{12.9}{eV}$。
\textbf{23.} O：原子 $6 \times (-13.6 \times 4.55^2/4) = \qty{-422.3}{eV}$；\ce{O+}，$Z^* =
4.90$，$5 \times (-13.6 \times 4.90^2/4) = \qty{-408.2}{eV}$；$I = \qty{14.2}{eV}$。模型值
11.5、12.9、14.2；实测值为 11.26、14.53、13.62。
\textbf{24.} 氧的实测能量低于氮：在氧中，
第四个 2p 电子必须与另一个电子共享一个轨道，二者之间的排斥
使其更容易被移除；而氮的三个未成对 p 电子
构成稳定的半充满亚层。斯莱特规则将同组内的所有电子
一视同仁，不考虑成对效应。
\textbf{25.} 氢原子的 1s 电子距离超过 $2a_0$ 的概率为
$\boldsymbol{13\mathrm e^{-4} \approx 0.238}$。
\end{solution}
